\chapter{System Design and Implementation}
\label{ch:method}

This chapter describes the design and implementation of the decentralized project management system proposed in this thesis. We first present the overall architecture (\Cref{sec:architecture}), then describe each layer in turn: the chaincode that governs on-chain state (\Cref{sec:chaincode-design}), the integration with IPFS for off-chain file storage (\Cref{sec:ipfs-integration}), the backend API that bridges client applications and the Fabric network (\Cref{sec:backend}), and the frontend interface used to interact with the system (\Cref{sec:frontend}). The architectural choices follow directly from the requirements identified in \Cref{sec:pm} and the technologies introduced in \Cref{ch:background}, and mirror the on-chain/off-chain separation validated by \textcite{adel2023} (\Cref{sec:adel}), adapted to the software project management data model.

% ─────────────────────────────────────────────
\section{System Architecture}
\label{sec:architecture}

The system follows a layered architecture in which each layer has a single, well-defined responsibility:

\begin{enumerate}
    \item \textbf{Frontend} — a browser-based client (\Cref{sec:frontend}) through which users create and inspect projects and tasks.
    \item \textbf{Backend API} — a stateless HTTP service (\Cref{sec:backend}) that translates REST requests into Fabric transactions and IPFS operations.
    \item \textbf{Hyperledger Fabric network} — a permissioned blockchain network of two organizations, each running one peer, and one ordering service, connected by a single channel. This network hosts the chaincode (\Cref{sec:chaincode-design}) that defines and enforces all business logic.
    \item \textbf{IPFS node} — a local content-addressed storage node (\Cref{sec:ipfs-integration}) used to store file attachments outside the ledger.
\end{enumerate}

The request flow for a typical write operation — for example, creating a task — proceeds as follows: the frontend sends an HTTP request to the backend; the backend constructs and submits a Fabric transaction using the Fabric Gateway client SDK; the transaction is endorsed by the peers of both organizations according to the channel's endorsement policy, ordered by the orderer, and committed to the ledger on both peers; the backend returns the resulting state to the frontend as a JSON response. Read operations follow the same path but use Fabric's \texttt{evaluateTransaction} call, which queries the local peer's world state directly without going through the ordering service, since no ledger update is required.

File attachment follows a variant of this flow: the file is first uploaded directly from the frontend to the local IPFS node, which returns a CID; this CID, together with the file name, is then submitted to the chaincode as part of a normal transaction, so that only the CID — not the file content — is recorded on the ledger (\Cref{sec:ipfs-integration}).

% ─────────────────────────────────────────────
\section{Chaincode Design}
\label{sec:chaincode-design}

The chaincode is implemented in JavaScript using the Fabric Shim API and deployed to a channel shared by both organizations in the network. It defines two primary record types, distinguished by a \texttt{docType} field: \emph{projects} and \emph{tasks}. Each record is stored under its own key in Fabric's world state, and every write to a key is simultaneously appended to that key's immutable transaction history, which Fabric retains independently of the current state (\Cref{sec:fabric}).

\subsection{Project Functions}

A project record holds an identifier, name, description, owner, a list of member identifiers, and a status (\texttt{active} or \texttt{archived}). The chaincode exposes:

\begin{itemize}
    \item \texttt{createProject(projectId, name, description, ownerId)} — creates a new project with the owner as its sole initial member, rejecting the call if the identifier is already in use.
    \item \texttt{getProject(projectId)} — returns the current state of a project.
    \item \texttt{addProjectMember(projectId, memberId)} — appends a new member to an active project, rejecting duplicates.
    \item \texttt{archiveProject(projectId)} — marks a project as archived; once archived, a project can no longer accept new tasks (enforced in \texttt{createTask}).
    \item \texttt{getProjectHistory(projectId)} — returns the full sequence of past states of a project, retrieved via Fabric's \texttt{getHistoryForKey} API, in reverse chronological order.
\end{itemize}

\subsection{Task Functions}

A task record holds an identifier, a reference to its parent project, a title, description, priority (\texttt{low}, \texttt{medium}, \texttt{high}), status, an optional assignee, and a list of file attachments. The chaincode exposes:

\begin{itemize}
    \item \texttt{createTask(taskId, projectId, title, description, priority)} — creates a task within an existing, active project.
    \item \texttt{getTask(taskId)} — returns the current state of a task.
    \item \texttt{assignTask(taskId, assigneeId)} — assigns a task to a project member; the chaincode verifies that the assignee is listed among the parent project's members before accepting the assignment, and rejects reassignment of a completed task.
    \item \texttt{updateTaskStatus(taskId, newStatus)} — transitions a task between statuses. Transitions are restricted to a fixed set of valid moves (\texttt{todo}~$\rightarrow$~\texttt{in-progress}, \texttt{in-progress}~$\rightarrow$~\texttt{todo} or \texttt{done}), so that, for instance, a task cannot move directly from \texttt{todo} to \texttt{done} or be reopened once \texttt{done}.
    \item \texttt{updateTaskMeta(taskId, field, value)} — allows editing of title, description, or priority on tasks that are not yet completed.
    \item \texttt{attachFile(taskId, fileName, ipfsCid)} — records a reference to an off-chain file (\Cref{sec:ipfs-integration}).
    \item \texttt{getTaskHistory(taskId)} — returns the full immutable history of a task, analogous to \texttt{getProjectHistory}.
\end{itemize}

\subsection{Determinism Constraints}

An important implementation constraint of chaincode is that its execution must be deterministic: every endorsing peer executes the same transaction independently and must produce byte-identical results for the transaction to be validated (\Cref{sec:fabric}). An early version of the chaincode included a server-generated timestamp (\texttt{new Date().toISOString()}) in each record. Because peers execute the same chaincode invocation at slightly different wall-clock times, this produced different output on different peers, causing endorsement mismatches and transaction failures. The timestamp fields were removed from the chaincode logic; where timestamps are required by the system, they are obtained from Fabric's own block metadata rather than generated inside the chaincode, which sidesteps the determinism issue.

% ─────────────────────────────────────────────
\section{IPFS Integration}
\label{sec:ipfs-integration}

Following the pattern established by \textcite{adel2023} (\Cref{sec:adel}), file content is never stored on the Fabric ledger. Instead, files are uploaded directly from the client to a local IPFS node via its HTTP API, which returns a content identifier (CID) computed as a cryptographic hash of the file's content (\Cref{sec:ipfs}). Only this CID, together with the original file name, is submitted to the chaincode's \texttt{attachFile} function and recorded as part of the task's on-chain state.

This design has two consequences that are relevant to the guarantees the system provides. First, because the CID is a function of file content, any subsequent modification of the file would produce a different CID; the CID recorded on the immutable ledger therefore constitutes a permanent, verifiable commitment to the exact content of the file at the time of attachment. Second, storing only a small fixed-size identifier on-chain, rather than the file itself, avoids the scalability and cost problems that would result from storing arbitrary binary data directly in blockchain state.

Retrieval works in the reverse direction: given a task record fetched from the chain, its attachments' CIDs can be used to fetch the corresponding files from any IPFS node holding a copy, including the local node's HTTP gateway, without any additional interaction with the Fabric network.

% ─────────────────────────────────────────────
\section{Backend API}
\label{sec:backend}

The backend is a stateless HTTP service that exposes the chaincode's functionality as a conventional REST API, so that client applications do not need to interact with the Fabric Gateway SDK directly. Each incoming request establishes a gRPC connection to a peer using the Fabric Gateway client, authenticated with the invoking identity's X.509 certificate and private key, submits or evaluates the corresponding chaincode transaction, and returns the result as JSON.

The API exposes endpoints mirroring the chaincode's function surface, for example \texttt{POST /projects} for \texttt{createProject}, \texttt{GET /tasks/\{id\}} for \texttt{getTask}, \texttt{PUT /tasks/\{id\}/status} for \texttt{updateTaskStatus}, and \texttt{GET /tasks/\{id\}/history} for \texttt{getTaskHistory}. Write operations map to Fabric's \texttt{submitTransaction} call, and read-only operations map to \texttt{evaluateTransaction}, preserving the distinction between operations that require network-wide ordering and consensus and those that can be answered from local peer state alone (\Cref{sec:fabric}).

This layer deliberately contains no business logic of its own: all validation (e.g., rejecting an invalid status transition, or an assignment to a non-member) is enforced by the chaincode, not the backend, ensuring that the guarantees of the system do not depend on trusting the backend implementation — a client could, in principle, bypass the backend and submit transactions directly to the network with equivalent results.

% ─────────────────────────────────────────────
\section{Frontend Interface}
\label{sec:frontend}

The frontend is a single-page application that provides a graphical interface for the operations described above: creating and loading projects, creating tasks within a loaded project, transitioning a task's status, attaching files, and inspecting a task's complete on-chain history.

The history view is the interface's central feature from the perspective of this thesis' contribution: rather than showing only the current state of a task, it renders every past state retrieved via \texttt{getTaskHistory}, in reverse chronological order, each annotated with its status and metadata at that point in time. This view makes directly visible, to a non-technical user, the property that motivated the system's design — that the record of how a task reached its current state cannot be altered after the fact, in contrast to the centralized platforms discussed in \Cref{sec:traditional-pm}.
